\documentclass[11pt,a4paper]{article}
\usepackage[italian]{babel}\usepackage[T1]{fontenc}\usepackage[utf8]{inputenc}
\usepackage{lmodern,amsmath,amssymb,tikz,enumitem,fancyhdr}
\usepackage[margin=1.7cm,headheight=14pt]{geometry}
\usetikzlibrary{arrows.meta}
\definecolor{rosso}{RGB}{210,25,30}
\definecolor{blu}{RGB}{28,91,119}
\definecolor{verde}{RGB}{35,128,80}
\pagestyle{fancy}\fancyhf{}\lhead{Matematica | 1SP}\rhead{Relazioni}\cfoot{\thepage}
\setlength{\parindent}{0pt}\setlength{\parskip}{5pt}
\setlist[enumerate]{label=\alph*),leftmargin=6mm,itemsep=3pt,topsep=4pt}
\newcommand{\ex}[2]{\vspace{9pt}{\large\bfseries\textcolor{rosso}{#1.}\ #2}\par}
\tikzset{v/.style={circle,fill=blu,inner sep=1.8pt},arr/.style={draw=rosso,-{Stealth[length=2.2mm]},thick}}
\newcommand{\faccia}[2]{%
\begin{scope}[xshift=#1cm]
\ifnum#2=1 \fill (0,0) circle (1.6pt);\fi
\ifnum#2=2 \fill (-.19,.19) circle (1.6pt);\fill (.19,-.19) circle (1.6pt);\fi
\ifnum#2=3 \fill (-.19,.19) circle (1.6pt);\fill (0,0) circle (1.6pt);\fill (.19,-.19) circle (1.6pt);\fi
\ifnum#2=4 \foreach \x in {-.19,.19}\foreach \y in {-.19,.19}\fill (\x,\y) circle (1.6pt);\fi
\end{scope}}
\newcommand{\tessera}[2]{\begin{tikzpicture}[baseline=-3mm]
\draw[rounded corners=2pt,thick,fill=white] (-.37,-.36) rectangle (1.11,.36);
\draw (.37,-.32)--(.37,.32);\faccia{0}{#1}\faccia{.74}{#2}
\end{tikzpicture}\hspace{2mm}}
\newcommand{\tesserina}[2]{\scalebox{.62}{\tessera{#1}{#2}}}
\newcommand{\cappioalto}[1]{\draw[arr] (#1.center) arc[start angle=-90,end angle=270,radius=.30];}
\newcommand{\cappiobasso}[1]{\draw[arr] (#1.center) arc[start angle=90,end angle=450,radius=.30];}
\begin{document}
\begin{center}{\LARGE\bfseries\textcolor{blu}{Relazioni e percorsi nei grafi}}\\[4pt] Relazioni riflessive e simmetriche, ponti e domino\end{center}
\textbf{Consegna generale.} Svolgi sul quaderno. Quando studi una proprietà, motiva la risposta: se non vale, indica un elemento o una coppia che lo dimostri.

\ex{1}{Coppie ordinate, prodotto cartesiano e grafo}
Considera $A=\{1,2,3\}$ e la relazione su $A$:
\[
R=\{(1,1),(1,2),(2,1),(2,2),(3,3)\}.
\]
\begin{enumerate}
\item Scrivi tutte le coppie di $A\times A$ ed evidenzia quelle che appartengono a $R$.
\item Rappresenta $R$ con un grafo orientato: la coppia $(x,y)$ corrisponde alla freccia da $x$ a $y$. Usa il rosso per le frecce.
\item La relazione è riflessiva? È simmetrica? Motiva.
\end{enumerate}

\ex{2}{Che cosa raccontano le frecce?}
Ciascun grafo rappresenta una relazione $R$ sull'insieme $A=\{a,b,c,d\}$. Studia i due casi separatamente.
\begin{center}\begin{tikzpicture}
\node at (1.6,3.25) {\textbf{Grafo I}};
\node[v,label=left:$a$] (a) at (0,2) {};\node[v,label=right:$b$] (b) at (3.2,2) {};\node[v,label=right:$c$] (c) at (3.2,0) {};\node[v,label=left:$d$] (d) at (0,0) {};
\path[arr] (a) edge (b) (b) edge[bend left=16] (c) (c) edge[bend left=16] (b);
\foreach \n in {a,b}{\cappioalto{\n}}
\foreach \n in {c,d}{\cappiobasso{\n}}
\begin{scope}[xshift=8.5cm]
\node at (1.6,3.25) {\textbf{Grafo II}};
\node[v,label=left:$a$] (e) at (0,2) {};\node[v,label=right:$b$] (f) at (3.2,2) {};\node[v,label=right:$c$] (g) at (3.2,0) {};\node[v,label=left:$d$] (h) at (0,0) {};
\path[arr] (e) edge[bend left=16] (f) (f) edge[bend left=16] (e) (g) edge[bend left=16] (h) (h) edge[bend left=16] (g);
\cappioalto{e}\cappiobasso{g}
\end{scope}
\end{tikzpicture}\end{center}
\begin{enumerate}
\item Scrivi $R$ come insieme di coppie ordinate, per ciascun grafo.
\item Stabilisci se $R$ è riflessiva e se è simmetrica.
\item Quali frecce devi aggiungere, al minimo, perché $R$ diventi sia riflessiva sia simmetrica? Elenca le coppie corrispondenti e conta le aggiunte.
\end{enumerate}

\ex{3}{Nella stessa classe}
Considera $A=\{\text{Anna},\text{Luca},\text{Marta},\text{Omar}\}$. Anna e Luca frequentano la 1SP, Marta la 1BS e Omar la 1AU. La relazione $R$ è definita da:
\[
x\,R\,y\quad\text{se e solo se $x$ frequenta la stessa classe di $y$.}
\]
\begin{enumerate}
\item Scrivi tutte le coppie di $R$ e disegna il grafo orientato.
\item $R$ è riflessiva? È simmetrica? Spiega usando il significato della relazione.
\item «Omar non frequenta la classe di nessun altro elemento di $A$: quindi dal suo vertice non parte alcuna freccia». Questa affermazione è corretta? Motiva.
\end{enumerate}
\newpage
\ex{4}{Relazioni in famiglia}
Paolo e Sara sono i genitori di Luca ed Emma. Luca ed Emma sono fratello e sorella. Paolo e Sara non sono fratelli; non ci sono altri rapporti di genitorialità o di fratellanza tra le quattro persone.

Lavora sull'insieme $A=\{\text{Paolo},\text{Sara},\text{Luca},\text{Emma}\}$. In ognuno dei due casi seguenti chiama la relazione $R$ e studiala separatamente.

\textbf{Caso I: essere fratello o sorella.}
\[
x\,R\,y\quad\text{se e solo se $x$ è fratello o sorella di $y$.}
\]
\begin{enumerate}
\item Elenca le coppie di $R$ e disegna il grafo, includendo tutti i vertici di $A$.
\item $R$ è riflessiva? È simmetrica? Motiva facendo riferimento alle persone.
\end{enumerate}

\textbf{Caso II: essere genitore.}
\[
x\,R\,y\quad\text{se e solo se $x$ è genitore di $y$.}
\]
\begin{enumerate}[resume]
\item Elenca le coppie di $R$ e disegna il grafo. Fai attenzione al verso delle frecce.
\item $R$ è riflessiva? È simmetrica? Motiva facendo riferimento alle persone.
\item Se aggiungi le frecce inverse a tutte quelle presenti, ottieni una relazione simmetrica? Ottieni anche una relazione riflessiva? Spiega entrambe le risposte.
\end{enumerate}

\ex{5}{Dalla regola al piano cartesiano}
Considera $A=\{-1,0,1\}$. La relazione $R$ su $A$ è definita da:
\[
x\,R\,y\quad\text{se e solo se }x+y=0.
\]
\begin{enumerate}
\item Scrivi tutte le coppie di $A\times A$ e seleziona quelle di $R$.
\item Rappresenta $R$ nel piano cartesiano qui sotto: per ogni coppia $(x,y)\in R$, segna un punto rosso. La prima coordinata va sull'asse orizzontale, la seconda su quello verticale.
\item Disegna anche il grafo orientato di $R$, usando gli elementi di $A$ come vertici. Spiega che cosa rappresenta una coppia ordinata nelle due rappresentazioni.
\item $R$ è riflessiva? È simmetrica? Motiva.
\item Quali coppie devi aggiungere, al minimo, per rendere $R$ riflessiva? La relazione ottenuta rimane simmetrica? Quante coppie contiene in tutto?
\end{enumerate}
\begin{center}\begin{tikzpicture}[scale=1.2]
\draw[step=.5,cyan!20,very thin] (-1.5,-1.5) grid (1.5,1.5);
\draw[-{Stealth},semithick] (-1.65,0)--(1.8,0) node[right] {$x$};
\draw[-{Stealth},semithick] (0,-1.65)--(0,1.8) node[above] {$y$};
\foreach \x in {-1,1}{\draw (\x,.05)--(\x,-.05);\node[below=3pt] at (\x,0) {$\x$};}
\foreach \y in {-1,1}{\draw (.05,\y)--(-.05,\y);\node[left=3pt] at (0,\y) {$\y$};}
\node[below left=2pt] at (0,0) {$0$};
\end{tikzpicture}\end{center}
\newpage
\begingroup\small
\setlength{\parskip}{3pt}
\setlist[enumerate]{label=\alph*),leftmargin=6mm,itemsep=2pt,topsep=3pt}
\begin{center}{\Large\bfseries\textcolor{blu}{Ponti e domino: si può fare il percorso?}}\end{center}
In questi problemi i lati non hanno un verso assegnato. Puoi visitare più volte un vertice, ma devi percorrere \textbf{ogni lato una e una sola volta}. Quando trovi un percorso, indica partenza e arrivo e segna il verso con frecce rosse.

\ex{6}{Sette ponti, come a K\"onigsberg}
La mappa schematica mostra quattro zone di terra e sette ponti numerati. All'interno di ciascuna zona puoi muoverti liberamente. Si può fare una passeggiata attraversando ogni ponte esattamente una volta?
\begin{center}\begin{tikzpicture}[scale=.95]
\fill[cyan!12] (0,0) rectangle (8,3.1);
\draw[verde,fill=green!5] (0,2.4) rectangle (8,3.1);
\draw[verde,fill=green!5] (0,0) rectangle (8,.6);
\draw[verde,fill=green!5] (2.7,1.5) ellipse (1.2 and .42);
\draw[verde,fill=green!5] (6.2,1.5) ellipse (.9 and .42);
\node[blu] at (4,2.75) {Zona $A$};\node[blu] at (4,.27) {Zona $C$};
\node[blu] at (2.7,1.5) {$B$};\node[blu] at (6.2,1.5) {$D$};
\tikzset{ponte/.style={draw=rosso,line width=4pt,line cap=rect}}
\draw[ponte] (2.1,2.4)--(2.1,1.86) node[midway,left=3pt,black] {\small 1};
\draw[ponte] (3.3,2.4)--(3.3,1.86) node[midway,right=3pt,black] {\small 2};
\draw[ponte] (2.1,.6)--(2.1,1.14) node[midway,left=3pt,black] {\small 3};
\draw[ponte] (3.3,.6)--(3.3,1.14) node[midway,right=3pt,black] {\small 4};
\draw[ponte] (6.2,2.4)--(6.2,1.92) node[midway,right=3pt,black] {\small 5};
\draw[ponte] (6.2,.6)--(6.2,1.08) node[midway,right=3pt,black] {\small 6};
\draw[ponte] (3.9,1.5)--(5.3,1.5) node[midway,above=3pt,black] {\small 7};
\end{tikzpicture}\end{center}
\begin{enumerate}
\item Trasforma la mappa in un grafo: che cosa rappresentano i vertici? E i lati? Numera i lati come i ponti e calcola il grado di ogni vertice.
\item Rispondi alla domanda iniziale usando il numero di vertici di grado dispari.
\item Viene chiuso il ponte 7. È possibile percorrere tutti i sei ponti rimasti una sola volta? Se sì, scrivi una sequenza di zone e di numeri dei ponti.
\item Riparti dalla mappa iniziale: tra quali due zone puoi costruire un nuovo ponte per rendere possibile la passeggiata? Scegli una possibilità e giustifica.
\end{enumerate}

\ex{7}{Domino: una catena e un anello}
Usa tutte le tredici tessere mostrate, una sola volta ciascuna. Puoi ruotarle: due estremità a contatto devono avere lo stesso numero di pallini.
\begin{center}
\tessera{0}{1}\tessera{1}{2}\tessera{2}{0}\tessera{0}{3}\tessera{3}{4}\tessera{4}{0}\tessera{2}{2}\\[5pt]
\tessera{1}{3}\tessera{1}{4}\tessera{2}{3}\tessera{2}{4}\tessera{0}{0}\tessera{4}{4}
\end{center}
\begin{enumerate}
\item Costruisci un grafo con un vertice per ciascun numero. Che cosa rappresenta ogni tessera? Come rappresenti le tessere doppie, per esempio \tesserina{2}{2}? Calcola tutti i gradi: un cappio contribuisce 2 al grado.
\item Puoi formare una catena con tutte le tessere? Puoi anche chiuderla ad anello? Giustifica con i gradi e scrivi un ordine possibile delle tessere.
\item Togli la tessera \tesserina{4}{0}. Puoi usare tutte le altre in una sola catena? Da quale numero devi partire e a quale devi arrivare, a meno di invertire il percorso? Scrivi una catena possibile.
\end{enumerate}

\ex{8}{Domino: basta tentare?}
\begin{center}
\tessera{0}{1}\tessera{0}{2}\tessera{0}{3}\tessera{1}{2}\tessera{1}{3}\tessera{2}{3}
\end{center}
\begin{enumerate}
\item Senza provare tutti gli ordini, stabilisci se puoi formare una sola catena con tutte le tessere. Disegna il grafo e usa i gradi per motivare.
\item Elimina una sola tessera a tua scelta: ora è possibile? Indica la tessera eliminata, i nuovi gradi e una catena che usi tutte le cinque tessere rimaste.
\item Riparti dalle sei tessere iniziali e aggiungi solo il doppio 2 \tesserina{2}{2}: il problema si risolve? Spiega che cosa cambia nei gradi.
\end{enumerate}
\endgroup
\end{document}
